\section{Cheapest Flights Within K Stops}
\label{sec:graph-cheapest-flights-k-stops}

\problemstatement{Given $n$ cities, a list of flights $[from, to,
price]$, a source, a destination, and an integer $k$, find the
cheapest price to travel from source to destination using \textbf{at
most $k$ stops} (i.e.\ at most $k{+}1$ edges). Return $-1$ if no such
route exists.}

\begin{patternbox}
\emph{``Shortest path, but bounded by the number of edges used''}
specifically \textbf{breaks} Dijkstra's greedy frontier-expansion
property: Dijkstra always finalizes whichever node has the globally
smallest distance, with no regard for how many edges that path used --
so it can permanently lock in a cheap-but-too-many-stops path, missing
a more-expensive-but-within-budget alternative. The fix: relax
\textbf{every} edge, for exactly $k{+}1$ \textbf{rounds}, copying
distances between rounds so each round adds at most one new edge to
every path -- this is a bounded version of Bellman-Ford's relaxation
strategy (Section~\ref{sec:graph-bellman-ford}).
\end{patternbox}

\subsection*{Why the copy between rounds is essential}

If edges were relaxed directly into the live $\textit{dist}$ array
within a single round, a single round could chain multiple
relaxations together (relax $u\to v$, then immediately use the updated
$\textit{dist}[v]$ to relax $v \to w$ in the \emph{same} round) --
silently allowing paths with more edges than the round count intends.
Using a fresh copy ($\textit{temp}$) as the round's write target, while
reading from the previous round's frozen $\textit{dist}$, guarantees
each round contributes \textbf{exactly one} additional edge to every
path under consideration.

\subsection*{Algorithm}

\begin{enumerate}
  \item Initialize $\textit{dist}[\textit{src}]=0$, all others
        $\infty$.
  \item Repeat $k{+}1$ times: copy $\textit{dist}$ into
        $\textit{temp}$; for every flight $(u,v,w)$, if
        $\textit{dist}[u] + w < \textit{temp}[v]$, update
        $\textit{temp}[v]$; then set $\textit{dist} \gets \textit{temp}$.
  \item Return $\textit{dist}[\textit{dst}]$ if finite, else $-1$.
\end{enumerate}

\subsection*{C++ implementation}

\begin{lstlisting}
#include <bits/stdc++.h>
using namespace std;

int findCheapestPrice(int n, vector<vector<int>>& flights, int src, int dst, int k) {
    vector<int> dist(n, INT_MAX);
    dist[src] = 0;

    for (int round = 0; round <= k; round++) {
        vector<int> temp = dist;
        for (auto& f : flights) {
            int u = f[0], v = f[1], w = f[2];
            if (dist[u] != INT_MAX && dist[u] + w < temp[v]) {
                temp[v] = dist[u] + w;
            }
        }
        dist = temp;
    }
    return dist[dst] == INT_MAX ? -1 : dist[dst];
}

int main() {
    vector<vector<int>> flights = {{0,1,100},{1,2,100},{0,2,500}};
    cout << findCheapestPrice(3, flights, 0, 2, 1) << endl; // 200
    return 0;
}
\end{lstlisting}

\subsection*{Dry run}

\dryrun{Take $\textit{flights}=[[0,1,100],[1,2,100],[0,2,500]]$,
$\textit{src}=0$, $\textit{dst}=2$, $k=1$ (at most $1$ stop, i.e.\ at
most $2$ edges).}

$\textit{dist}=[0,\infty,\infty]$. Round $0$ ($k{+}1=2$ rounds total):
$\textit{temp}=[0,\infty,\infty]$; relax $0{\to}1$:
$0+100=100<\infty$, $\textit{temp}[1]=100$; relax $1{\to}2$:
$\textit{dist}[1]=\infty$ still (reading the \emph{frozen} pre-round
value), skip; relax $0{\to}2$: $0+500=500<\infty$,
$\textit{temp}[2]=500$. $\textit{dist}=[0,100,500]$. Round $1$:
$\textit{temp}=[0,100,500]$; relax $0{\to}1$: $100$ not $<100$, no
change; relax $1{\to}2$: $\textit{dist}[1]+100=100+100=200<500$,
$\textit{temp}[2]=200$; relax $0{\to}2$: $500$ not $<200$, no change.
$\textit{dist}=[0,100,200]$. Final answer: $\textit{dist}[2]=200$ --
the path $0{\to}1{\to}2$ (one stop, total $200$), correctly preferred
over the direct $0{\to}2$ edge costing $500$.

\begin{complexitybox}
\textbf{Time Complexity.} $O(k \cdot E)$ -- $k{+}1$ rounds, each
scanning every flight once.
\textbf{Space Complexity.} $O(V)$ for \textit{dist} and \textit{temp}.
\end{complexitybox}

\begin{takeawaybox}
This problem is a useful warning against reaching for Dijkstra
reflexively on every ``cheapest path'' phrasing: the moment a
constraint couples the answer to \emph{how the path is structured}
(here, edge count) rather than purely to its total weight, Dijkstra's
greedy finalization can be provably wrong, and a bounded, round-based
relaxation (Bellman-Ford's core idea, capped early) becomes the correct
tool.
\end{takeawaybox}
